% META: {"id": "TNSI-HISTOIRE-INFORMATIQUE-QCM", "chapitre": "TNSI-HISTOIRE-INFORMATIQUE", "type_objet": "qcm", "genere_depuis": "chapitres/TNSI-HISTOIRE-INFORMATIQUE/qcm/TNSI-HISTOIRE-INFORMATIQUE-QCM.json", "status": "generated"}
% Fichier genere par scripts/build_qcm_tex.py — ne pas editer a la main.

\section*{\textcolor{chapcolor}{\MakeUppercase{Histoire de l'informatique — QCM d'auto-évaluation}}}

\begin{center}
\textit{Pour chaque question, une seule réponse est exacte.}
\end{center}

\begin{enumerate}

\item[] \textbf{Capacité C1}

\item \textbf{[Q1]} Alan Turing formalise le concept de machine universelle en :
  \begin{enumerate}[label=\Alph*.]
    \item 1936.
    \item 1948.
    \item 1969.
    \item 1991.
  \end{enumerate}

\item \textbf{[Q2]} ARPANET, ancêtre d'Internet, est mis en service en :
  \begin{enumerate}[label=\Alph*.]
    \item 1948.
    \item 1969.
    \item 1989.
    \item 2000.
  \end{enumerate}

\item[] \textbf{Capacité C2}

\item \textbf{[Q3]} Depuis les débuts de l'informatique, le rôle du matériel a évolué vers :
  \begin{enumerate}[label=\Alph*.]
    \item une spécialisation croissante, chaque machine ne traitant qu'une seule tâche.
    \item une disparition progressive au profit du seul logiciel.
    \item davantage de généricité, la logique particulière étant portée par le logiciel.
    \item aucune évolution notable.
  \end{enumerate}

\item \textbf{[Q4]} Internet et le World Wide Web sont dans le rapport suivant :
  \begin{enumerate}[label=\Alph*.]
    \item ce sont deux noms de la même chose.
    \item ce sont deux réseaux physiques distincts et indépendants.
    \item Internet est un service qui fonctionne au-dessus du Web.
    \item le Web est un service qui fonctionne au-dessus du réseau Internet.
  \end{enumerate}

\end{enumerate}

% NEXUS-QCM-TEACHER-ONLY-BEGIN
\ifnxVersionProfesseur
\clearpage
\section*{Clé de correction — réservée au professeur}

\begin{center}
\begin{tabular}{lll}
\hline
Question & Capacité & Réponse exacte \\
\hline
Q1 & C1 & \textbf{A} \\
Q2 & C1 & \textbf{B} \\
Q3 & C2 & \textbf{C} \\
Q4 & C2 & \textbf{D} \\
\hline
\end{tabular}
\end{center}

\subsection*{Diagnostic des réponses erronées}

\begin{itemize}
  \item \textbf{Q1}
  \begin{itemize}
    \item \textbf{B} — En 1948, le Manchester Baby exécute un programme enregistré en mémoire. Le modèle théorique de Turing le précède de plus de dix ans. \emph{Renvoi : C1, repères chronologiques.}
    \item \textbf{C} — 1969 est la mise en service d'ARPANET, un jalon des réseaux et non du modèle de calcul. \emph{Renvoi : C1, repères chronologiques.}
    \item \textbf{D} — 1991 est l'année de la première diffusion publique de Python, créé en 1989, et de la diffusion du Web hors du CERN : bien après le modèle théorique de Turing. \emph{Renvoi : C1, repères chronologiques.}
  \end{itemize}
  \item \textbf{Q2}
  \begin{itemize}
    \item \textbf{A} — 1948 correspond à la première exécution du Manchester Baby, ordinateur électronique à programme enregistré ; ARPANET est mis en service en 1969. \emph{Renvoi : C1, repères chronologiques.}
    \item \textbf{C} — 1989 est l'année où Tim Berners-Lee propose le Web au CERN. ARPANET, l'un des précurseurs d'Internet, a été mis en service vingt ans auparavant. \emph{Renvoi : C1, Internet et le Web.}
    \item \textbf{D} — ARPANET est mis en service en 1969 ; l'année 2000 se situe 31 ans après ce repère. \emph{Renvoi : C1, repères chronologiques.}
  \end{itemize}
  \item \textbf{Q3}
  \begin{itemize}
    \item \textbf{A} — C'est l'inverse du mouvement historique : les premières machines étaient câblées pour une tâche, et le programme enregistré les a rendues polyvalentes. \emph{Renvoi : C2, rôles relatifs du logiciel et du matériel.}
    \item \textbf{B} — Tout logiciel s'exécute sur du matériel. Sa part relative dans la valeur a crû, mais le support n'a pas disparu. \emph{Renvoi : C2, rôles relatifs du logiciel et du matériel.}
    \item \textbf{D} — Le déplacement de la logique du câblage vers le logiciel est l'une des évolutions majeures de la discipline. \emph{Renvoi : C2, rôles relatifs du logiciel et du matériel.}
  \end{itemize}
  \item \textbf{Q4}
  \begin{itemize}
    \item \textbf{A} — Le courriel et bien d'autres services empruntent Internet sans passer par le Web : les deux ne se confondent donc pas. \emph{Renvoi : C2, Internet et le Web.}
    \item \textbf{B} — Le Web n'est pas un réseau physique : c'est un ensemble de pages liées, échangées grâce au protocole HTTP. \emph{Renvoi : C2, Internet et le Web.}
    \item \textbf{C} — L'élève intervertit les deux niveaux. L'infrastructure est Internet ; le Web est l'un des services qu'elle transporte. \emph{Renvoi : C2, Internet et le Web.}
  \end{itemize}
\end{itemize}
\fi
% NEXUS-QCM-TEACHER-ONLY-END
